\section{When to Evolve}
\label{sec:when}
The temporal dimension of self-evolution in LLM-based agents mainly concerns the relationship between learning processes and task execution. Therefore, the second key aspect of a self-evolving agent is identifying the \textit{evolving timing}, i.e., at which stage the self-evolving strategy $f$ is invoked and applied to the agent system. To this end, we propose a taxonomy that distinguishes between two temporal modes of self-evolution: Intra-test-time self-evolution and inter-test-time self-evolution.

\textbf{Intra-test-time self-evolution} refers to adaptive processes that occur during task execution, where agents recognize their limitations on a specific problem and initiate targeted learning mechanisms to enhance their capabilities in real-time \citep{xi2024enhancing, bi2024forest}. This mode of evolution is characterized by its immediate coupling with the task at hand: the agent improves its problem-solving abilities for a specific problem encountered, creating a dynamic interplay between performance and adaptation.

\textbf{Inter-test-time self-evolution} refers to learning processes that occur between task completions, leveraging accumulated experiences to improve future performance. This category encompasses diverse methodological approaches: offline learning paradigms that extract knowledge from pre-collected datasets through iterative refinement \citep{zelikman2022star, zelikman2024quiet}, and online learning paradigms that continuously adapt based on streaming interaction data \citep{qi2024webrl, qiu2025alita, qian2024new, wang2025mobile}. 

The implementation of self-evolution across these temporal phases leverages three fundamental learning paradigms in LLMs: in-context learning (ICL) \citep{dong2022survey, min2021metaicl, wies2023learnability}, which adapts behavior through contextual examples without modifying parameters; supervised fine-tuning (SFT), which updates model weights through gradient-based optimization on labeled data \citep{devlin2018bert, shen2024rethinking, dong2023abilities}; and reinforcement learning (RL), which shapes behavior through reward-driven policy optimization \citep{kaelbling1996reinforcement, sun2024llm,zhang2025nemotron}. While these learning paradigms remain conceptually consistent across temporal contexts, their instantiation differs in terms of data availability and learning objectives:


Intra-test-time is characterized by its online nature: learning data emerges dynamically during task execution, with optimization directly targeting performance enhancement on the immediate problem instance. This real-time coupling necessitates rapid adaptation mechanisms that can process learning data and feedback signals and modify behavior within the temporal constraints of active task-solving. On the other hand, inter-test-time is characterized by its retrospective nature: learning algorithms operate on historical data, whether from curated datasets or accumulated behavioral trajectories, with optimization objectives oriented toward improving expected performance across the task distribution rather than maximizing success on any specific problem instance. This temporal decoupling enables more sophisticated learning procedures that can identify cross-task patterns, consolidate diverse experiences, and develop generalizable capabilities without the immediacy constraints of active task execution.

\begin{figure}[t]
  \centering
  \includegraphics[width=0.8\linewidth]{Figures/when.pdf}
  \caption{An overview of when to evolve. The top pathway illustrates intra-test-time self-evolution, where adaptation (e.g., variant generation, verification, and policy update) occurs within task execution. The bottom pathway depicts inter-test-time self-evolution, where learning happens retrospectively through rollout, trajectory analysis, and policy updates. }
  \label{fig:when}
\end{figure}

\subsection{Intra-Test-Time Self-Evolution}

In intra-test-time self-evolution, agents engage in self-improvement processes that are intrinsically coupled with solving the immediate task at hand. The distinguishing characteristic of this temporal phase is its synchronous nature: feedback signals are generated and processed during task execution, with optimization objectives specifically targeted at improving performance on the current problem instance rather than generalizing to future tasks. Here, we introduce how the three learning paradigms are realized in this temporal phase.

\paragraph{In-Context Learning} Intra-test-time ICL methods leverage the model's context window as a dynamic memory system for immediate adaptation without parameter modification. These approaches typically employ self-reflective mechanisms where agents analyze their own performance, generate verbal critiques or insights, and maintain these reflections in episodic memory buffers to guide subsequent decisions within the same task context \citep{shinn2023reflexion, madaan2023self_refine}. Some methods extend beyond simple reflection to include dynamic planning revision, where agents can modify their entire approach based on environmental feedback, switching between action execution and plan modification as needed. For instance, AdaPlanner \citep{sun2023adaplanner} decomposes tasks into manageable sub-goals and predicts environmental feedback for each. During execution, its refiner component distinguishes between in-plan feedback (observations aligning with predictions) and out-of-plan feedback (deviating observations). For in-plan feedback, the refiner dynamically queries the LLM through a specialized \texttt{ask\_LLM()} action to parse observations and extract pertinent information. For out-of-plan feedback, the refiner proactively revises the entire plan and resumes solving from an intermediate point, rather than restarting from scratch. This adaptive closed-loop framework eliminates the need for prior knowledge about feedback structures and enables more efficient decision-making. Similarly, TrustAgent \citep{hua2024trustagent} employs rule-based plan revision during execution, modifying its approach based on language feedback to evolve toward safer planning strategies. These ICL methods demonstrate how test-time adaptation can achieve sophisticated behavioral modification without permanent model changes, maintaining flexibility while preserving the model's general capabilities.

\paragraph{Supervised Fine-Tuning.} Intra-test-time SFT represents a paradigm shift where models perform immediate self-modification through learned meta-adaptation strategies. Self-adaptive language modeling \citep{zweiger2025self} exemplifies this approach by generating ``self-edits'', which are meta-level instructions that can restructure information representations, specify optimization hyperparameters, or invoke tools for data augmentation and gradient computation. These self-edits trigger immediate supervised fine-tuning, resulting in persistent weight updates that adapt the model to the current task. The key innovation lies in the meta-learning phase, where reinforcement learning trains models to produce effective self-edits by using the downstream performance of the updated model as the reward signal, essentially teaching models how to teach themselves. \citet{acikgoz2025self} introduce a Test-Time Self-Improvement (TT-SI) framework that enables agents to adapt on-the-fly by first identifying uncertain test samples through self-awareness. For these challenging inputs, the agent then generates a single synthetic training example and performs a temporary, lightweight parameter update to improve its immediate performance before resetting its weights.


\paragraph{Reinforcement Learning.} Intra-test-time RL enables models to develop new capabilities on-demand when encountering problems beyond their current competence. LADDER \citep{simonds2025ladder} demonstrates this through its test-time reinforcement learning (TTRL) mechanism: upon identifying a particularly challenging problem, the system generates a focused set of related problem variants and conducts intensive, targeted reinforcement learning specifically for that problem class. This approach transforms insurmountable challenges into learning opportunities, allowing models to expand their problem-solving repertoire during deployment rather than failing or providing suboptimal solutions. The method represents a form of just-in-time skill acquisition, where computational resources are invested precisely when and where they are needed most.


\subsection{Inter-Test-Time Self-Evolution}

Inter-test-time self-evolution represents the predominant learning process in autonomous agents, wherein adaptation occurs following task execution rather than during it. In this temporal mode, agents complete a given task, extract feedback signals, including explicit rewards \citep{gao2024designing}, gradients \citep{amari1993backpropagation, bottou2010large}, and performance metrics \citep{ge2023openagi}, and subsequently leverage this information to enhance their capabilities for future problem-solving. This retrospective learning process decouples task performance from capability improvement, allowing agents to consolidate experiences, identify patterns of success and failure, and systematically refine their behavioral policies without the computational constraints imposed by real-time task demands.

\paragraph{In-Context Learning.} Inter-test-time in-context learning has emerged as a widely adopted approach for agent self-improvement. This paradigm leverages execution results and feedback from previous tasks as contextual information for future problem-solving. Wang et al. \citep{wang2024agent} demonstrate this principle by inducing workflows from agent action histories and incorporating them into the context for subsequent tasks. The field of in-context reinforcement learning (ICRL) \citep{moeini2025survey, laskin2022context, lee2023supervised} extends this concept by maintaining histories of observations and actions within the agent's context window. These methods exploit the hypothesis that pre-trained neural networks can implement implicit reinforcement learning algorithms within their forward pass, processing contextual information to adapt behavior without parameter updates \citep{kirsch2023towards}. A defining characteristic of ICRL is in-context improvement: the phenomenon whereby agent performance progressively enhances as task-relevant information accumulates in the context, enabling sophisticated adaptation through attention mechanisms rather than gradient-based learning.


\paragraph{Supervised Fine-Tuning.} Inter-test-time SFT~\citep{chen2025sft} methods establish a paradigm of iterative self-improvement through synthetic data generation and self-evaluation. SELF \citep{lu2023self} pioneered meta-cognitive training, where models first acquire self-feedback and self-refinement capabilities, then iteratively generate responses to unlabeled instructions and enhance them through self-critique. STaR \citep{zelikman2022star} and Quiet-STaR \citep{zelikman2024quiet} focus on reasoning improvement through rationalization—models attempt problems, then generate explanations for correct answers they initially failed to solve, creating augmented training data that combines successful attempts with post-hoc reasoning. SiriuS \citep{zhao2025sirius} extends this to sequential problem-solving, maintaining repositories of correct solutions while augmenting failures through multi-stage refinement involving feedback incorporation, regeneration, and rephrasing. These methods share a core insight: models can bootstrap their own improvement by learning to evaluate and enhance their outputs, creating high-quality training signals from initially imperfect attempts without extensive human supervision. Recent frameworks such as ARIA \citep{he-etal-2025-enabling} further extend this paradigm by incorporating human-in-the-loop guidance into test-time adaptation, allowing agents to proactively identify knowledge gaps and request expert feedback.


\paragraph{Reinforcement Learning.} Inter-test-time RL leverages unconstrained computational resources to optimize agents through extensive environmental interaction and sophisticated curriculum design. RAGEN \citep{wang2025ragen} and DYSTIL~\citep{wang2025dystildynamicstrategyinduction}  employ online reinforcement learning for multi-turn interactive tasks, continuously refining policies through on-policy learning in simulated dialogues. Learning Like Humans \citep{zhang2025learning} introduces cognitive-inspired training with adaptive difficulty progression, combining on-policy exploration with off-policy efficiency and expert demonstrations to accelerate learning. Domain-specific applications demonstrate the versatility of inter-test-time RL: WebRL \citep{qi2024webrl} develops web navigation agents through self-evolving curricula that automatically adjust task complexity based on performance, while DigiRL \citep{bai2024digirl} enables device-control agents to master in-the-wild interactions through autonomous reinforcement learning. 
These approaches exploit the pre-deployment phase to engage in extensive trial-and-error learning, developing robust policies through thousands of interactions that would be impractical during real-time deployment.
